\documentclass[../../main.tex]{subfiles}
\begin{document}

{\bf [解答]}

曲線$C$ を
$
\begin{pmatrix} a \\ b
\end{pmatrix} \neq
\begin{pmatrix} 0 \\ 0
\end{pmatrix}
$
だけ平行移動した図形は，
\begin{align}
  y-b = (x-a)^3 - (x-a) \quad (-1+a \le x \le 1+a) \label{eq:1}
\end{align}
である．

これと $y=x^3-x$ との共有点がただ1つある時，まず $y$ を消去して，
\begin{align}
  x^3 - x - b &= (x-a)^3 - (x-a) \\
  -x - b &= -3a x^2 + (3a^2-1) x - a^3 + a \\
  \therefore
  b &= 3a x^2 - 3a^2 x + a^3 - a \equiv f(x) \label{eq:2}
\end{align}
である．この二次方程式$b=f(x)$が$-1 \le x \le 1$ かつ $-1+a \le x \le 1+a$ にただ1つ解を持つ条件を調べる．下図のように，$f(x)$の軸$x = \dfrac{a}{2}$ は常にこの範囲にある．

\begin{figure}[htb]
  \centering
  \begin{tikzpicture}[scale=1.9, >=stealth, every node/.style={font=\small}]

    % --- 1. 座標軸 ---
    \draw[->, thick] (-2.5, 0) -- (2.6, 0) node[right] {$a$};
    \draw[->, thick] (0, -1.5) -- (0, 1.5) node[above] {$x$};
    \node[below left=2pt] at (0,0) {$O$};

    % --- 2. 共通領域の平行四辺形 ---
    % 境界線: x = 1, x = -1, x = a-1, x = a+1
    % 頂点: (-2, -1), (0, -1), (2, 1), (0, 1)
    \coordinate (V1) at (-2, -1);
    \coordinate (V2) at (0, -1);
    \coordinate (V3) at (2, 1);
    \coordinate (V4) at (0, 1);

    % 斜線ハッチング塗りつぶし
    \fill[pattern=north east lines, pattern color=blue!65, opacity=0.75]
    (V1) -- (V2) -- (V3) -- (V4) -- cycle;

    % 領域の境界線
    \draw[thick, gray!70] (-2.3, 1) -- (2.3, 1) node[right, font=\scriptsize] {$x = 1$};
    \draw[thick, gray!70] (-2.3, -1) -- (2.3, -1) node[right, font=\scriptsize] {$x = -1$};
    \draw[thick, blue!80!black] (-2.3, -1.3) -- (0.5, 1.5) node[above right, font=\scriptsize] {$x = a+1$};
    \draw[thick, blue!80!black] (-0.5, -1.5) -- (2.3, 1.3) node[below right, font=\scriptsize] {$x = a-1$};

    % --- 3. 放物線の軸直線 x = a/2 ---
    \draw[line width=1.6pt, red!85!black] (-2.2, -1.1) -- (2.2, 1.1);
    \node[red!90!black, font=\footnotesize, above left] at (1.5, 0.75) {$x = \frac{a}{2}$};

    % --- 4. 補助線と目盛り ---
    \draw[dashed, gray!60] (-2, 0) -- (-2, -1) -- (0, -1);
    \draw[dashed, gray!60] (2, 0) -- (2, 1) -- (0, 1);

    \node[below=2pt, font=\scriptsize] at (-2, 0) {$-2$};
    \node[below=2pt, font=\scriptsize] at (-1, 0) {$-1$};
    \node[below=2pt, font=\scriptsize] at (1, 0) {$1$};
    \node[below=2pt, font=\scriptsize] at (2, 0) {$2$};

    \node[left=2pt, font=\scriptsize] at (0, 1) {$1$};
    \node[left=2pt, font=\scriptsize] at (0, -1) {$-1$};

  \end{tikzpicture}
  \caption{共通定義域（斜線部）と放物線 $f(x)$ の軸 $x = \dfrac{a}{2}$ の関係}
  \label{fig:common_domain_axis}
\end{figure}

よって唯一つ解を持つ条件は$b$が$f(x)$の頂点に等しい時で
\begin{align}
  f\Bigl(\dfrac{a}{2}\Bigr) = b \iff b = \dfrac{1}{4}a^3 - a \label{eq:3}
\end{align}
である．
この時 $(a,b) \ne (0,0)$ だから $a \ne 0$ である．
\cref{eq:1}に代入して，
\begin{align}
  y = (x-a)^3 - (x-a) + \dfrac{1}{4}a^3 - a \quad (-1+a \le x \le 1+a) \label{eq:4}
\end{align}
である．
$P(X, Y)$ とおくと\cref{eq:4}が $P$ を通ることから，
\begin{align}
  Y = (X-a)^3 - (X-a) + \dfrac{1}{4}a^3 - a \quad (-1+a \le X \le a+1) \label{eq:5}
\end{align}
である．
これをみたす $a(\ne 0)$ がただ3つ存在する条件を考える．
\cref{eq:5}の右辺を $g(a)$ とおいて$a$についての三次方程式とみなして考える．
\begin{align}
  g(a) &= -\dfrac{3}{4}a^3 + 3X a^2 - 3X^2 a + X^3 - X \\
  g'(a) &= -\dfrac{9}{4}a^2 + 6X a - 3X^2 \\
  &= -\dfrac{3}{4}(3a-2X)(a-2X)
\end{align}
である．\cref{eq:5}が$a$について$3$つの解を持つには，$g'(a)=0$なる点$a=\dfrac{2X}{3},2X$が区間内$(-1+a \le X \le a+1)$に存在しないといけない．
これらを図示すると下図である．

\begin{figure}[htb]
  \centering
  \begin{tikzpicture}[scale=1.2, >=stealth, every node/.style={font=\small}]

    % --- 1. 座標軸 ---
    \draw[->, thick] (-3.3, 0) -- (3.3, 0) node[right] {$X$};
    \draw[->, thick] (0, -2.8) -- (0, 2.8) node[above] {$a$};
    \node[below left=2pt] at (0,0) {$O$};

    % --- 2. 帯領域 X-1 <= a <= X+1 ---
    % 帯のハッチング塗りつぶし
    \fill[pattern=north east lines, pattern color=blue!65, opacity=0.75]
    (-3.2, -2.2) -- (-3.2, -4.2) -- (3.2, 2.2) -- (3.2, 4.2) -- cycle;

    % 境界線
    \draw[very thick, blue!85!black, domain=-3.2:1.7]
    plot (\x, {\x + 1}) node[above right, font=\scriptsize] {$a = X+1$};
    \draw[very thick, blue!85!black, domain=-1.7:3.2]
    plot (\x, {\x - 1}) node[below right, font=\scriptsize] {$a = X-1$};

    % --- 3. 停留点の直線 ---
    % a = 2X
    \draw[line width=1.5pt, red!85!black, domain=-1.35:1.35]
    plot (\x, {2*\x}) node[above, font=\footnotesize] {$a = 2X$};

    % a = (2/3)X
    \draw[line width=1.5pt, teal!85!black, domain=-3.0:3.0]
    plot (\x, {2.0/3.0*\x}) node[above right, font=\footnotesize] {$a = \frac{2}{3}X$};

    % --- 4. 交点・端点のプロットと補助破線 ---
    % a = 2X と a = X+1 の交点: (1, 2)
    % a = 2X と a = X-1 の交点: (-1, -2)
    \draw[dashed, gray!60] (1, 0) -- (1, 2) -- (0, 2);
    \draw[dashed, gray!60] (-1, 0) -- (-1, -2) -- (0, -2);

    % a = (2/3)X と a = X-1 の交点: (3, 2)
    % a = (2/3)X と a = X+1 の交点: (-3, -2)
    \draw[dashed, gray!60] (3, 0) -- (3, 2);
    \draw[dashed, gray!60] (-3, 0) -- (-3, -2);

    % 目盛り
    \node[below=2pt, font=\scriptsize] at (1, 0) {$1$};
    \node[above=2pt, font=\scriptsize] at (-1, 0) {$-1$};
    \node[below=2pt, font=\scriptsize] at (3, 0) {$3$};
    \node[above=2pt, font=\scriptsize] at (-3, 0) {$-3$};
    \node[left=2pt, font=\scriptsize]  at (0, 2) {$2$};
    \node[right=2pt, font=\scriptsize] at (0, -2) {$-2$};

    % 端点の白抜き丸 (開区間)
    \node[circle, draw, thick, fill=white, red!85!black, inner sep=1.2pt] at (1, 2) {};
    \node[circle, draw, thick, fill=white, red!85!black, inner sep=1.2pt] at (-1, -2) {};
    \node[circle, draw, thick, fill=white, teal!85!black, inner sep=1.2pt] at (3, 2) {};
    \node[circle, draw, thick, fill=white, teal!85!black, inner sep=1.2pt] at (-3, -2) {};

  \end{tikzpicture}
  \caption{帯領域 $X-1 \le a \le X+1$ と停留点直線 $a=2X, a=\dfrac{2}{3}X$ の関係}
  \label{fig:stationary_points_in_strip}
\end{figure}

従って，図から$-1 < X < 1 \ (X \ne 0)$ が必要で，$X$ の正負で場合分けする．

\subparagraph{$1^\circ \ -1 < X < 0$ の時}

この時の$g(a)$の増減表は以下のようになる．

\begin{table}[htb]
  \centering
  \begin{tabular}{c|c|c|c|c|c|c|c}
    $a$ & $X-1$ & $\dots$ & $2X$ & $\dots$ & $\dfrac{2}{3}X$ & $\dots$ & $X+1$ \\ \hline
    $g'$ & & $-$ & $0$ & $+$ & $0$ & $-$ & \\ \hline
    $g$ & & $\searrow$ & 極小 & $\nearrow$ & 極大 & $\searrow$ &
  \end{tabular}
  \caption{$1^\circ$ における $g$ の増減表}
\end{table}

したがって，\cref{eq:5}が $a \ne 0$ に3実解を持つ条件は，
\begin{align}
  \begin{cases}
    \min\Bigl(g(2X), g(X+1)\Bigr) \le Y \le \max\Bigl(g(X-1), g\Bigl(\dfrac{2}{3}X\Bigr)\Bigr) \\
    Y \ne g(0) \\
    Y \ne g(2X) \\
    Y \ne g(X+1)
  \end{cases}\label{eq:6}
\end{align}
である．

\paragraph{$2^\circ \ 0 < X < 1$ の時}

この時の$g(a)$の増減表は以下のようになる．

\begin{table}[htb]
  \centering
  \begin{tabular}{c|c|c|c|c|c|c|c}
    $a$ & $X-1$ & $\dots$ & $\dfrac{2}{3}X$ & $\dots$ & $2X$ & $\dots$ & $X+1$ \\ \hline
    $g'$ & & $-$ & $0$ & $+$ & $0$ & $-$ & \\ \hline
    $g$ & & $\searrow$ & 極小 & $\nearrow$ & 極大 & $\searrow$ &
  \end{tabular}
  \caption{$2^\circ$ における $g$ の増減表}
\end{table}

したがって，\cref{eq:5}が $a \ne 0$ に3実解を持つ条件は
\begin{align}
  \begin{cases}
    \min\Bigl(g\Bigl(\dfrac{2}{3}X\Bigr), g(X+1)\Bigr) \le Y \le \max\Bigl(g(2X), g(X-1)\Bigr) \\
    Y \ne g(0) \\
    Y \ne g(2X) \\
    Y \ne g\Bigl(\dfrac{2}{3}X\Bigr)
  \end{cases} \label{eq:7}
\end{align}
である．
\casesend

$g$の値は以下のようになる．
\begin{align}
  g(X+1) &= \dfrac{1}{4}(X+1)^3 - (X+1) + (-1)^3 - (-1) \\
  &= \dfrac{1}{4}(X+1)(X+3)(X-1) \label{eq:8} \\
  g(X-1) &= \dfrac{1}{4}(X-1)^3 - (X-1) \\
  &= \dfrac{1}{4}(X-1)(X+1)(X-3) \label{eq:9} \\
  g(2X) &= \dfrac{1}{4} \cdot 8 X^3 - 2X + (-X)^3 - (-X) \\
  &= X^3 - X \label{eq:10} \\
  g\Bigl(\dfrac{2}{3}X\Bigr) &= \dfrac{1}{4}\Bigl(\dfrac{2}{3}X\Bigr)^3 - \Bigl(\dfrac{2}{3}X\Bigr) + \Bigl(\dfrac{1}{3}X\Bigr)^3 - \Bigl(\dfrac{1}{3}X\Bigr) \\
  &= \dfrac{1}{9}X^3 - X \label{eq:11} \\
  g(0) &= X^3-X
\end{align}
従って，もとめる領域は下図斜線部 (境界は実線のみ含む)である．

\begin{figure}[htb]
  \centering
  \begin{tikzpicture}[scale=3.5, >=stealth, every node/.style={font=\small}]

    % --- 1. 座標軸と目盛り ---
    \draw[->, thick] (-1.6, 0) -- (1.7, 0) node[right] {$x$};
    \draw[->, thick] (0, -1.6) -- (0, 1.6) node[above] {$y$};
    \node[below left=2pt] at (0, 0) {$O$};

    \node[below=3pt, font=\scriptsize] at (-1, 0) {$-1$};
    \node[below=3pt, font=\scriptsize] at (1, 0) {$1$};

    % y 軸の -1/4 の目盛り (手書きスケッチの「-1/4」を反映)
    \draw[gray!70] (-0.05, -0.25) -- (0.05, -0.25);
    \node[left=2pt, font=\scriptsize, gray!90!black] at (0, -0.25) {$-\frac{1}{4}$};

    % --- 2. 求める領域のハッチング塗りつぶし (手書き image_8.png の影部分を忠実に再現) ---
    % 領域 1: -1 < x < 0 の範囲
    %   上端: y = g(x-1) = 0.25*(x-1)*(x+1)*(x-3)  (境界②)
    %   下端: y = max( g(2x), g(0) ) = x^3 - x     (境界①)  および y = g(2x/3) = (1/9)x^3 - x (境界④)
    %   ※ -0.577 < x < 0 では y = (1/9)x^3 - x が上に来るため、そこが下端境界。
    %   ※ -1 < x < -0.577 では y = x^3 - x が下端境界。
    \fill[pattern=north east lines, pattern color=blue!65, opacity=0.8]
    plot[domain=-1:-0.57735, samples=20] (\x, {0.25*(\x-1)*(\x+1)*(\x-3)})
    -- plot[domain=-0.57735:0, samples=20] (\x, {(1.0/9.0)*\x*\x*\x - \x})
    -- plot[domain=0:-1, samples=20] (\x, {\x*\x*\x - \x})
    -- cycle;

    \fill[pattern=north east lines, pattern color=blue!65, opacity=0.8]
    (0, 0)
    -- plot[domain=0:0.57735, samples=20] (\x, {(1.0/9.0)*\x*\x*\x - \x})
    -- plot[domain=0.57735:1, samples=20] (\x, {0.25*(\x-1)*(\x+1)*(\x+3)})
    -- plot[domain=1:0, samples=20] (\x, {\x*\x*\x - \x})
    -- cycle;

    % --- 3. 境界曲線のプロットと詳細な数式ラベル ---
    % ①: y = x^3 - x (実線)
    \draw[very thick, blue!85!black, domain=-1.25:1.25, samples=70]
    plot (\x, {\x*\x*\x - \x});
    \node[blue!95!black, font=\scriptsize, right=3pt] at (1.1, {1.1^3 - 1.1}) {$y = x^3 - x$};

    % ②: y = 0.25(x-1)(x+1)(x-3) (破線)
    \draw[thick, red!85!black, dashed, domain=-1.25:1.15, samples=60]
    plot (\x, {0.25*(\x-1)*(\x+1)*(\x-3)});
    \node[red!95!black, font=\scriptsize, below=4pt] at (0.9, 1.0) {$y = \frac{1}{4}(x-1)(x+1)(x-3)$};

    % ③: y = 0.25(x+1)(x+3)(x-1) (点線/二点鎖線のように細かい破線)
    \draw[thick, teal!85!black, dotted, domain=-1.15:1.25, samples=60]
    plot (\x, {0.25*(\x+1)*(\x+3)*(\x-1)});
    \node[teal!95!black, font=\scriptsize, above right=1pt] at (-1.05, -1) {$y = \frac{1}{4}(x+1)(x+3)(x-1)$};

    % ④: y = (1/9)x^3 - x (一点鎖線状の破線)
    \draw[thick, gray!70, dash pattern=on 4pt off 3pt, domain=-1.35:1.35, samples=60]
    plot (\x, {(1.0/9.0)*\x*\x*\x - \x});
    \node[gray!90!black, font=\scriptsize, above left=1pt] at (1.2, {(1.0/9.0)*1.2^3 - 1.2}) {$y = \frac{1}{9}x^3 - x$};

    % --- 4. 領域に含まれない除外点 (白抜き丸: draw=black, fill=white で完全に白抜き) ---
    % 横軸上の端点・原点: (-1,0), (0,0), (1,0)
    \node[circle, draw, thick, black, fill=white, inner sep=1pt] at (-1, 0) {};
    \node[circle, draw, thick, black, fill=white, inner sep=1pt] at (0, 0) {};
    \node[circle, draw, thick, black, fill=white, inner sep=1pt] at (1, 0) {};

    % 左側の極大周辺の主要交点 (x = -1/sqrt(3) ≈ -0.577, y = 0.25*(-1.577)*(-0.577)*(-3.577)? )
    % 交点として曲線②と④の交点に該当する点
    \coordinate (P_ex1) at (-0.57735, {0.25*(-1.57735)*(0.42265)*(-3.57735)});
    \node[circle, draw, thick, black, fill=white, inner sep=1pt] at (P_ex1) {};

    % 右側の極小周辺の主要交点 (x = 1/sqrt(3) ≈ 0.577)
    % 曲線①と③の交点に該当する点
    \coordinate (P_ex2) at (0.57735, {0.25*(1.57735)*(3.57735)*(-0.42265)});
    \node[circle, draw, thick, black, fill=white, inner sep=1pt] at (P_ex2) {};

  \end{tikzpicture}
  \caption{条件を満たす点 $P(X,Y)$ の存在領域（斜線部）．境界は実線部のみ含み白抜きの点は含まない．}
  \label{fig:xy_swept_region_final}
\end{figure}

{\bf [別解]}

前半の\cref{eq:3}を別の方法で導く．

$\vec{A} =
\begin{pmatrix} a \\ b
\end{pmatrix}$ だけ平行移動するとする．
$OA$ の中点 $M\Bigl(\dfrac{1}{2}a, \dfrac{1}{2}b\Bigr)$ とする．
$C$ の任意の点 $(x,y)$ の，$M$ に関する対称点を $(x',y')$ とすると，
\begin{align}
  x = a - x', \quad y = b - y'
\end{align}
だから，$y = x^3 - x$ に代入して
\begin{align}
  y - b = (x'-a)^3 - (x'-a)
\end{align}
である．
つまり $(x',y')$ は $C'$ 上の点である．
したがって，$C, C'$ は $M$ を中心として点対称である．
よって，$C, C'$ の共有点も $M$ を中心として点対称だから，
\begin{align}
  \text{「共有点がただ1つ} \iff \text{共有点が} M\text{」}
\end{align}
となり，
\begin{align}
  \dfrac{b}{2} = \left(\dfrac{a}{2}\right)^3-\dfrac{a}{2} \quad \therefore b = \dfrac{1}{4}a^3 - a
\end{align}
を得る．

\end{document}
